\chapter{Proof overview}\label{sec:overview}

In this section, we explain the main steps of the proofs and the ideas behind them. We do not restate the theorems. For each step, we point to the result where it is proved. In \cref{sec:ov-slowness}, we compare the two kinds of slowness that the two halves of the paper address. We explain the proofs of \cref{thm:main-condensed,thm:main-initial} in \cref{sec:ov-condensed} and the proofs of \cref{thm:main-memory,thm:main-plateau,thm:main-cusp} in \cref{sec:ov-prestatic}; these two parts can be read in either order. \Cref{sec:ov-dependencies} lists the dependencies among the main propositions.

Both proofs have the same structure. We first show, by a short deterministic argument, that a few properties of the Gibbs measure, or of the critical points of the field, imply the dynamical conclusion. We then prove these properties in separate sections, with tools from the statics of spherical spin glasses. In the condensed phase, the deterministic argument is the reduction of \cref{sec:reduction}, and most of the work goes into its inputs. Below the static threshold, the probabilistic input is a set of estimates on the constrained pressure (\cref{sec:pressure}), and most of the new arguments are deterministic (\cref{sec:spectral}).

We write $\approx$ in the displays that are heuristics. Such a display indicates the size of a quantity to leading exponential order, or the origin of a formula; it is not used in the proofs. All other displays are exact statements, which we prove where indicated.

\section{Two kinds of slowness}\label{sec:ov-slowness}

We can measure the distance to equilibrium from a worst-case starting point, from the uniform law $\nu_N$, or from a starting point drawn from the Gibbs measure (\cref{sec:intro-starts}). The two halves of the paper concern the second and the third of these. We use different mechanisms for them.

We first explain why an exponentially small spectral gap does not by itself slow down the dynamics started from $\nu_N$. Let $f=\dd\nu_N/\dd\pi=Z_{N,\beta}e^{-\beta H_N}$. By reversibility, the density of $\nu_NP_t$ with respect to $\pi$ is $P_tf$. Let $(\phi_j)_{j\ge0}$ be an orthonormal basis of eigenfunctions of $-L$ in $L^2(\pi)$ with $\phi_0=1$ (\cref{lem:pre-dynamics}). Expanding $P_tf-1$ in this basis and using the Cauchy--Schwarz inequality, we find
\begin{equation}\label{eq:ov-L2}
\norm{\nu_NP_t-\pi}_{\TV}=\frac12\int\abs{P_tf-1}\dd\pi\le\frac12\Bigl(\sum_{j\ge1}e^{-2\lambda_jt}\ip f{\phi_j}^2\Bigr)^{1/2}.
\end{equation}
We see that a small eigenvalue $\lambda_j$ delays the relaxation from $\nu_N$ only through the coefficient $\ip f{\phi_j}$. Small coefficients on the slow modes remove that particular obstruction to relaxation; a fast-mixing bound also requires control of the remaining terms in \eqref{eq:ov-L2}. The example of \cite{GheissariGrantYu2026} recalled in \cref{sec:intro-starts} shows that a small spectral gap can coexist with fast mixing from the uniform law. To bound the distance from below, we need a different argument.

We use the variational form of the total variation distance. Since $\norm{\mu-\pi}_{\TV}=\sup_A(\pi(A)-\mu(A))$, the distance between $\nu_NP_t$ and $\pi$ is close to one if and only if there exists a set $A$ such that $\pi(A)$ is close to one and $(\nu_NP_t)(A)$ is close to zero. Thus separation at a given time is equivalent to small endpoint mass, for the dynamics started from $\nu_N$, on a set carrying almost all of the Gibbs measure. Note that this set must depend on the disorder: for a deterministic set $A$, rotation invariance shows that $\pi(A)$ and $(\nu_NP_t)(A)$ both have mean $\nu_N(A)$ over the disorder. By the contraction of total variation (\cref{lem:pre-contraction}), it suffices to exhibit such a set at the final time. In the condensed phase, we take for $A$ a union of thin bands around a bounded number of critical points of a field derived from $H_N$. In the proof, we establish the stronger property that the path avoids the caps around these critical points in each fixed window. Below the static threshold, the Gibbs measure is predicted to be spread over exponentially many states, each of exponentially small mass, and a set that carries nearly all of the Gibbs measure must meet exponentially many of them. Our conditioned hitting estimate has a polynomial error for each state, which cannot be summed over exponentially many states. This is one reason why our results on the dynamics started from $\nu_N$ are restricted to the condensed phase.

For Gibbs-typical starting points, the law of $X_t$ is $\pi$ at all times, and the memory of the dynamics lies in the joint law of $X_0$ and $X_t$. Here we use a slowly varying observable. If $B\colon S_N\to\R^N$ has small vector energy, then, for the stationary diffusion,
\begin{equation}\label{eq:ov-endpoint-energy}
\E_\pi\abs{B(X_t)-B(X_0)}^2=2\ip B{(1-P_t)B}\le2t\,\cE_{\rm vec}(B)
\end{equation}
(\cref{lem:sp-trap}). For every fixed $\rho>0$ and each time $t$, Markov's inequality implies that $\Pp_\pi\{\abs{B(X_t)-B(X_0)}>\rho\}\le2t\,\cE_{\rm vec}(B)/\rho^2$. The endpoints are close in probability when $t\,\cE_{\rm vec}(B)\to0$; this is not a bound on the supremum along a trajectory. If, in addition, $B(x)$ is aligned with the direction of $x$, in the sense that $U(x)\cdot B(x)\ge q>0$ for most $x$, then $B(X_t)$ records where the dynamics started. In \cref{part:prestatic}, we construct one such observable with exponentially small energy and derive all five conclusions of \cref{thm:main-memory} from it.

\section{The condensed phase}\label{sec:ov-condensed}

Let $(\xi,\beta)$ be strictly 1RSB with parameters $x,Q$ and a nondegenerate contact at $Q$ (\cref{def:intro-certificate}). Recall the center field $H_Q(u)=H_N(\sqrt Q\,u)$ and the radial derivative $K_Q(u)$, which is the derivative of $t\mapsto H_N(t\sqrt Q\,u)$ at $t=1$, for $u\in S_N$. The proof of \cref{thm:main-condensed} rests on a picture of the Gibbs measure that Subag established for pure models at sufficiently low temperature \cite{Subag2017geometry}. In this picture, the Gibbs measure is carried by the bands $B_u=\{x\in S_N:\abs{R(x,u)-\sqrt Q}\le b\}$ around the critical points $u$ of $H_Q$ whose values lie within a bounded distance of the maximum of $H_Q$. We call these critical points centers.

We can see the reason in the static results of \cref{sec:thermal}. By \cref{prop:th-structure}(5), we have
\begin{equation}\label{eq:ov-tap}
F(\beta)=\beta E_0+\frac12\log(1-Q)+\frac{\beta^2}2\psi_Q(1),
\end{equation}
where $E_0$ is the limit of $N^{-1}\max_{S_N}H_Q$ and $\psi_Q(t)=\xi(Q+(1-Q)t)-\xi(Q)-(1-Q)\xi'(Q)t$ is the residual covariance. Each term has a meaning for a single band. Conditionally on a critical point $u$ with $H_Q(u)=NE$, the field at the points of latitude $\sqrt Q$ around $u$ is $NE$ plus a residual Gaussian field with covariance $N\psi_Q$ (\cref{lem:mb-regression}). The coarea density at latitude $\sqrt Q$ contributes the entropy term $\log(1-Q)/2$ per site, and a latitude interval of width $N^{-1/2}$ around $\sqrt Q$ has Haar mass $(1-Q)^{N/2+o(N)}$. For the fixed-width band, the weighted latitude integral must also be controlled by the saddle estimates. The residual field is strictly replica symmetric (\cref{prop:th-structure}(2)), which suggests that its partition function is close to its mean $e^{N\beta^2\psi_Q(1)/2}$. Heuristically, we find
\begin{equation}\label{eq:ov-band-weight}
\int_{B_u}e^{\beta H_N}\dd\nu_N\approx e^{\beta H_Q(u)}(1-Q)^{N/2}e^{N\beta^2\psi_Q(1)/2}\approx e^{NF(\beta)}e^{\beta(H_Q(u)-NE_0)}.
\end{equation}
Comparing with \eqref{eq:ov-tap}, we see that a single band around a center of nearly maximal value carries a Gibbs weight comparable to the whole partition function, to leading exponential order. We control the number of centers with the Kac--Rice formula. For non-pure mixtures, \cref{thm:mo-microscopic} shows that the expected number of local maxima of $H_Q$, with normalized radial derivative in its admissible interval $U$ and value above $m_N+t$, tends to $e^{-c_0t}$ for each fixed $t$, where $c_0=\beta x$. Here $m_N$ is chosen in that theorem by requiring the expected number of these local maxima above $m_N$ to equal one, and $m_N=NE_0-(\log N)/(2c_0)+O(1)$. In this heuristic, the bands of the centers $u$ with energy deficit $m_N-H_Q(u)$ close to $v$ have total weight of order $e^{c_0v}e^{-\beta v}$, which is summable in $v$ because $x<1$. Note that the ratio $c_0/\beta=x$ is the mass of the atom at zero of the Parisi measure. We expect that a finite number of bands carries most of the Gibbs measure, and that the bands of centers with a larger energy deficit carry a part that decays exponentially in the deficit.

On the other hand, a point drawn from $\nu_N$ has overlap of order $N^{-1/2}$ with each center. Fix $a<\sqrt Q-b$, and let $C_a(u)=\{x:\abs{R(x,u)}\ge a\}$ be the union of the two caps of latitude $a$ around $\pm u$. Every band $B_u$ lies in $C_a(u)$, and the diffusion started from $\nu_N$ must enter $C_a(u)$ before it reaches $B_u$. The heart of the proof is that, before time $e^{\gamma N}$, it does not enter $C_a(u)$ for any center.

\subsection{The separation criterion}\label{sec:ov-separation}

For $M>0$, let $\cC_{N,M}$ be the set of critical points of $H_Q$ with value in $[m_N-M,m_N+M]$ and normalized radial derivative in a small interval $U$ around $K_0$, and let $B_{N,M}$ be the union of their bands (\cref{def:red-centers}). For a time $T_N$, we write $\theta_N(H,u)=\Pp^H_{\nu_N}\{X_t\in C_a(u)\text{ for some }t\le T_N\}$ for the probability that the diffusion started from $\nu_N$ enters $C_a(u)$ before $T_N$. Testing $B_{N,M}$ in the definition of the total variation distance and taking a union bound over centers, we have, for every realization of the field,
\begin{equation}\label{eq:ov-separation}
\inf_{0\le t\le T_N}\norm{\nu_NP_t-\pi}_{\TV}\ge1-\pi\bigl(B_{N,M}^c\bigr)-\sum_{u\in\cC_{N,M}}\theta_N(H_N,u)
\end{equation}
(\cref{prop:red-consumer}). It suffices to prove three properties. The first is that the bands carry almost all of the Gibbs mass, in the sense that $\E\pi(B_{N,M}^c)$ tends to zero as $N\to\infty$ and then $M\to\infty$. The second is that the expected number of centers in each fixed window is bounded. The third is that the expected sum in \eqref{eq:ov-separation} tends to zero for each fixed $M$.

The third property concerns an average over critical points of a random field. We compute it with the Kac--Rice formula, which expresses it through the law of the field conditioned to have a critical point at a fixed $u$ with prescribed marks, that is, prescribed value and radial derivative. Let $D_u=\abs{\det\nabla^2_{S_N}H_Q(u)}$, and let $\E_{h,k,u}$ denote the expectation with respect to the field conditioned on $H_Q(u)=Nh$, $K_Q(u)=Nk$, and $\nabla_{S_N}H_Q(u)=0$ (\cref{def:kr-conditioned}). Using \cref{thm:kr-marked} and rotation invariance, we can write this average as
\begin{equation}\label{eq:ov-kac-rice}
\E\sum_{u\in\cC_{N,M}}\theta_N(H_N,u)=c_N\int_{\mathcal W_{N,M}}p_N(h,k)\,\E_{h,k,u}\bigl[D_u\,\theta_N(H_N,u)\bigr]\dd h\dd k,
\end{equation}
where $\mathcal W_{N,M}$ is the set of normalized marks of the window, $p_N$ is their density at a point, and $c_N$ collects the volume of $S_N$ and the density of the gradient at zero. For pure models, we have $K_Q=pH_Q$, and the integral is over the value alone. We meet two difficulties here. The first is that the weight $\theta_N$ is the probability of a trajectory event. It is a bounded Borel function of the whole field. We do not assume continuity or dependence only on the value, radial derivative, gradient, and Hessian at $u$. The marked formulas cited in \cref{sec:kacrice} impose continuity or tameness hypotheses. We prove the required formula for Borel weights in the coefficient space of the polynomial field, with an explicit version of the conditional law (\cref{prop:kr-coefficient,thm:kr-marked}).

The second difficulty is the determinant. The conditioned hitting estimate, which we describe in \cref{sec:ov-entrance}, bounds $q_N=\sup_{h,k,u}\E_{h,k,u}\theta_N$, without the weight $D_u$. We cannot simply drop the weight, since $D_u$ and the trajectory depend on the same field. To overcome this difficulty, we apply H\"older's inequality to the normalized determinant. If the normalized radial derivative $k$ lies a fixed distance above the edge $2\sqrt{b_Q}$ of the semicircle law, where $b_Q=Q^2\xi''(Q)$, then every fixed moment of $D_u/\E_{h,k,u}D_u$ is bounded uniformly in $N$ (\cref{prop:kr-determinant}). For $r>1$, this implies that
\begin{equation}\label{eq:ov-holder}
\E_{h,k,u}\bigl[D_u\theta_N\bigr]\le\mathsf K_r^{1/r}q_N^{1-1/r}\,\E_{h,k,u}D_u
\end{equation}
(\cref{lem:kr-holder}). We insert \eqref{eq:ov-holder} into \eqref{eq:ov-kac-rice} and apply \eqref{eq:ov-kac-rice} once more with $\theta_N=1$. We find that the average is at most $\mathsf K_r^{1/r}q_N^{1-1/r}\E\abs{\cC_{N,M}}$ (\cref{thm:kr-transfer}). This is a form of the conditional Cauchy--Schwarz argument of Huang and Sellke \cite[Corollary~3.16]{HuangSellke2024}. Since $q_N\le CN^{-(d-2)/2}$, it is at most $C_{M,r}N^{-(d-2)(1-1/r)/2}$ (\cref{prop:cp-transfer}).

The order in which we choose the parameters is part of the argument. The concentration property lets $M$ tend to infinity after $N$. We must choose the time $T_N=e^{\gamma N}$ before $M$, and the bound on the expected sum must hold for every fixed $M$ with the same $\gamma$; a hitting estimate whose exponent depended on $M$ would not suffice. We prove the hitting estimate uniformly over a compact rectangle $\mathcal K=[E_0-1,E_0+1]\times U$ of normalized marks, which contains the marks of every fixed window once $N$ is large (\cref{cor:ent-rectangle}). We fix the interval $U$, the band width $b$, and the rectangle first, then the cap latitude $a$ and the exponent $\gamma$, and only then the tolerance and the window half-width $M$. Only the constants and the threshold for $N$ depend on $M$ (\cref{tab:cp-parameters}).

\subsection{Statics}\label{sec:ov-thermal}

We now explain where $Q$, $E_0$, $K_0$, and $c_0$ come from. Recall that the Crisanti--Sommers functional has a unique minimizer $\zeta$, the Parisi measure, and that $\zeta$ is characterized by its first variation: the function $V_\zeta$ of \cref{sec:intro-model} is nonnegative and vanishes on the support of $\zeta$ (\cref{lem:pre-support}). The 1RSB assumptions of \cref{def:intro-certificate} ask that the two-atom measure $x\delta_0+(1-x)\delta_Q$ pass this test, that its contact set $\{V=0\}$ be exactly $\{0,Q\}$, and that the contact at $Q$ be nondegenerate, $V''(Q)>0$. They concern one explicit function of one variable. We verify them for the classes of \cref{cor:main-classes} in \cref{sec:th-classes,app:certificate}.

The first variation also describes the geometry of the Gibbs measure. Following Subag \cite{Subag2024}, we say that an overlap $q$ is multisamplable if, for every fixed $k$, the probability that $k$ independent samples from $\pi$ have pairwise overlaps close to $q$ is not exponentially small. Let $V_\beta$ be the first-variation function of the Parisi measure, whose zero set is the contact set. In \cref{thm:th-contact}, we show for every finite mixture and every $\beta$ that $q$ is multisamplable if and only if $V_\beta(q)=0$, that is, if and only if $q$ belongs to the contact set. The proof uses Subag's generalized TAP functional $T_\beta(q)$, which is the sum of $\beta$ times the ground energy at radius $\sqrt q$, the entropy term $\log(1-q)/2$, and the free energy of the residual mixture $\psi_q$. Subag proved that $T_\beta(q)\le F(\beta)$, with equality exactly when $q$ is multisamplable \cite[Theorem~5]{Subag2024}. The ground energy and the residual free energy are the values of two variational problems. We show in \eqref{eq:th-telescope} that $F(\beta)$ is the same sum with both problems evaluated at trial points built from the Parisi measure, and this gives an upper bound for $T_\beta(q)$. If $V_\beta(q)>0$, then a variation of the ground profile in an explicit direction lowers the zero-temperature Crisanti--Sommers functional, and this leads to the bound $T_\beta(q)\le F(\beta)-V_\beta(q)^2/(8\beta^3M_q)$, with an explicit $M_q>0$ (\cref{thm:th-contact}(3)). If $V_\beta(q)=0$, we show that both trial points are optimal. That every zero of $V_\beta$ is multisamplable is due to Subag, who also proved the converse for generic mixtures \cite[Theorem~10]{Subag2024}.

We collect the static consequences of the 1RSB assumptions in \cref{prop:th-structure}. The multisamplable overlaps are $0$ and $Q$, and the residual model at $Q$ is strictly replica symmetric. The ground state of $H_Q$ is strictly one-step replica symmetry breaking at zero temperature, with energy $E_0$ and normalized radial derivative $K_0$ given by explicit formulas, and $c_0=\beta x<\beta$. Finally, $K_0$ lies strictly above the edge of the semicircle law: with $y=Q/(\beta(1-Q))$, we have
\begin{equation}\label{eq:ov-gap}
K_0-2\sqrt{b_Q}=\frac{(y-\sqrt{b_Q})^2}y>0.
\end{equation}
This gap matters for the following reason. Conditionally on $H_Q(u)=Nh$, $K_Q(u)=Nk$, and $\nabla_{S_N}H_Q(u)=0$, the Hessian $\nabla^2_{S_N}H_Q(u)$ is a scaled GOE matrix with spectrum near $[-2\sqrt{b_Q},2\sqrt{b_Q}]$, shifted by $-k$ (\cref{lem:th-hessian}). For a deterministic $u$ and prescribed marks with $k$ near $K_0$, the Hessian is bounded above by a fixed negative multiple of the identity with conditional probability $1-Ce^{-cN}$ (\cref{cor:kr-hessian}). The statements about all high critical points are established separately in \cref{prop:mo-all-high,lem:pc-hessian}. The normalized determinant moments used in \eqref{eq:ov-holder} follow from \cref{prop:kr-determinant}.

\subsection{The directing measure}\label{sec:ov-directing}

The concentration argument of \cref{sec:ov-capture} places a Gibbs sample in the band of a center by finding additional independent replicas whose overlaps with each other and with the given sample are all close to $Q$. Such replicas exist if the overlap structure of the Gibbs measure has no other scale. Recall the theorem of Dovbysh and Sudakov (\cref{thm:dir-DS}): every limit in distribution of the overlap array of independent Gibbs samples is generated by drawing a random measure $G$ on the unit ball of a Hilbert space, and then independent vectors $v_1,v_2,\dots$ with law $G$, and setting $R_{ij}=v_i\cdot v_j$ for $i\ne j$. We show in \cref{thm:dir-norm} that, under the 1RSB assumptions, $G$ is carried by the vectors of squared norm $Q$. The vectors in a small ball around $v_1$ then have mutual inner products close to $Q$, and they produce these replicas.

We have to exclude two possibilities. Directing vectors of a squared norm $q\in(0,1)\setminus\{Q\}$ would make $q$ multisamplable, and the exponential bound for replicas at overlaps outside the contact set rules them out (\cref{lem:th-replica-exclusion,lem:dir-exclusion}); we exclude vectors of norm one with an entropy estimate (\cref{lem:dir-cap}). An atom of $G$ at the origin would describe a positive fraction of samples that are asymptotically orthogonal to all other samples. To exclude it, we use a single identity of Ghirlanda--Guerra type, with the test function $\xi(R)$, which we derive from Gaussian integration by parts and the self-averaging of the energy (\cref{lem:dir-GG}). It states that, for every limiting overlap law, every $n\ge1$, and every bounded function $f$ of $(R_{ij})_{1\le i<j\le n}$,
\begin{equation}\label{eq:ov-GG}
\E\bigl[f\xi(R_{1,n+1})\bigr]=\frac{(1-x)\xi(Q)}n\E f+\frac1n\sum_{j=2}^n\E\bigl[f\xi(R_{1j})\bigr].
\end{equation}
Let $a_n$ be the probability that $\xi(R_{1j})=0$ for $2\le j\le n$, and take for $f$ the indicator of this event. The sum on the right vanishes. After the first step, we know that $\abs{R_{ij}}\le Q$ for $i\ne j$, which implies that $\xi(R_{1,n+1})\le\xi(Q)$, and the left side is at most $\xi(Q)(a_n-a_{n+1})$. We find that $a_{n+1}\le(1-(1-x)/n)a_n$, and hence $a_n\le n^{-(1-x)}\to0$. A sample with directing vector zero would satisfy $\xi(R_{1j})=0$ for all $j$. We conclude that $G$ has no atom at the origin. The argument uses no ultrametricity and no description of the atoms or the weights of $G$.

\subsection{Critical points near the maximum}\label{sec:ov-counts}

We prove the second property of \cref{sec:ov-separation}, bounded counts in fixed windows, in \cref{sec:purecover} for pure models and in \cref{sec:mixone,sec:mixpair} for non-pure mixtures, together with the existence of at least one center, which the concentration argument needs. The following limits for the complexity in two variables and for the normalized intensity concern non-pure mixtures. For pure models, $K_Q=pH_Q$, and \cref{prop:pc-counts,lem:pc-tight} provide the count bound in the energy variable alone and the tightness of the centered maximum. The first moment is a one-point Kac--Rice computation. Ben Arous, Subag, and Zeitouni computed the exponential rate of the expected number of critical points of $H_Q$ with normalized value $e$ and normalized radial derivative $k$, which is the marked complexity \cite{BenArousSubagZeitouni2020}. We show that, for each $e$, it is strongly concave in $k$, that its maximizer at $e=E_0$ is $K_0$, and that its derivative in $e$ at $(E_0,K_0)$ is $-c_0$. A Laplace computation with bounded prefactors then shows that, for all fixed $a<b$,
\begin{equation}\label{eq:ov-first-moment}
\E\#\bigl\{u:\nabla_{S_N}H_Q(u)=0,\ K_Q(u)/N\in U,\ H_Q(u)-m_N\in[a,b]\bigr\}\longrightarrow e^{-c_0a}-e^{-c_0b}
\end{equation}
(\cref{thm:mo-microscopic}). In particular, the expected count in every fixed window is bounded, and the probability that the maximum of $H_Q$ exceeds $m_N+L$ is at most $e^{-c_0L}$ in the limit.

For existence, we need a second moment. For pure models, existence follows from the tightness of the centered maximum, which is a consequence of the work of Subag and Zeitouni \cite{SubagZeitouni2017} (\cref{lem:pc-tight,prop:pc-counts}). For non-pure mixtures, we bound the second moment of the number $C_N(L)$ of critical points in the window $[m_N-L,m_N+L]$ that satisfy a selection condition described below. The two-point Kac--Rice integrand is singular as the overlap $r$ of the two points tends to $\pm1$. We split the pairs into three ranges.
\begin{itemize}
\item \emph{Overlaps $r\ge\eps$.} We use the exclusion of high-energy points at positive overlap, \cref{thm:mo-filter}, which is proved at the ground-state parameters $(E_0,K_0)$ and follows the conditional trial argument of Huang and Sellke \cite[Section~3.3]{HuangSellke2024}. We show that, conditionally on a critical point at $u$ with marks near $(E_0,K_0)$, with probability $1-e^{-cN}$ no point at overlap between $\eps$ and $1-\eta$ with $u$ has value larger than $NE_0$ minus a fixed multiple of $N$, and no other critical point has overlap at least $1-\eta$ with $u$. The local maxima with this property are called retained maxima (\cref{def:mo-retained}). Distinct retained maxima have overlap smaller than $\eps$, and the selection removes only an exponentially small fraction of the expected count, also after weighting by the determinant (\cref{prop:mo-filter-loss}).
\item \emph{Overlaps $r\le-\delta$.} For even mixtures, these pairs are reflections of pairs with $r\ge\delta$. In general, an elementary bound for spherical codes in a cap suffices \cite{Rankin1955,BargMusin2007}. Suppose that some unit vectors have pairwise inner products at most $\eps<\delta^2$, and that $v_1,\dots,v_k$ among them have inner product at most $-\delta$ with a fixed $v_0$. Then $k\delta\le\abs{v_1+\dots+v_k}\le(k+k(k-1)\eps)^{1/2}$, which implies that each retained maximum has at most $(1-\eps)/(\delta^2-\eps)$ such points, deterministically (\cref{lem:mp-gram}).
\item \emph{Overlaps $\abs r<\delta$.} We show that, for $\abs r\le N^{-\alpha}$ with $1/4<\alpha<1/2$, the two-point integrand is the product of the two one-point integrands and the density of the overlap of two independent uniform points, up to a factor $1+O(N^{1-4\alpha}+N^{-1/2})$ (\cref{thm:mp-local-pair}). On $\abs r\le N^{-\alpha}$, the overlap density satisfies $f_N(r)=(1+o(1))\sqrt{N/(2\pi)}\,e^{-Nr^2/2}$, uniformly in $r$. On $N^{-\alpha}\le\abs r<\delta$, \cref{thm:mp-local-pair}(c) bounds the full two-point integrand by $C\sqrt N\,e^{-Nr^2/4}$ times the product of the one-point integrands, using the overlap density together with the regression and determinant estimates. For the factorization, we need the expected product of the two Hessian determinants under two-point conditioning. We write each Hessian as a bulk block bordered by one row and a corner entry and compare the bulk determinants with a maximal correlation inequality for Gaussian vectors, whose constant does not depend on the dimension.
\end{itemize}
Combining the three ranges, we find that $\E C_N(L)^2\le(1+o(1))(\E C_N(L))^2+(1+K)\E C_N(L)$, with a constant $K$ that does not depend on $L$ (\cref{thm:mp-count}). Since $\E C_N(L)\to e^{c_0L}-e^{-c_0L}$, Chebyshev's inequality implies that the window contains a retained maximum with probability at least $1-(1+K)/(e^{c_0L}-e^{-c_0L})-o(1)$ (\cref{thm:mc-counts}).

\subsection{Gibbs concentration on bands}\label{sec:ov-capture}

We prove the concentration property in \cref{sec:purecover} for pure models (\cref{thm:pc-capture}) and in \cref{sec:mixbands} for non-pure mixtures (\cref{thm:mb-capture}). For pure models at sufficiently low temperature, this concentration was proved by Subag \cite[Theorem~1]{Subag2017geometry}; we need it at every temperature of the condensed phase and in a fixed window. For non-pure mixtures, the three inputs below verify \cref{prop:red-capture}. The pure proof instead starts from a coarse cap cover, lets the energy deficit per spin tend to zero, removes the caps around critical points with a large energy deficit, and passes from the caps of a fixed window to bands of width of order $N^{-1/2}$ (\cref{lem:pc-diagonal,sec:pc-capture}).

The first step is a coarse cover: for every fixed energy tolerance $\tau>0$, the caps (in pure models, \cref{prop:pc-coarse}) or the bands (in non-pure mixtures, \cref{prop:mb-coarse}) around the critical points of $H_Q$ with value at least $N(E_0-\tau)$ carry almost all of the Gibbs mass. By \cref{thm:dir-norm}, among a fixed number of additional independent replicas, a Gibbs sample typically has $n$ whose mutual overlaps and overlaps with it are close to $Q$. Using Subag's estimates for the free energy of replicated bands with a fixed number of replicas \cite[Lemmas~17--19 and Proposition~24]{Subag2024}, we show that the normalized empirical mean of such a group has value close to $NE_0$ under $H_Q$ (\cref{prop:pc-replica-means,lem:mb-replica}). Writing $\bar\sigma=n^{-1}\sum_i\sigma^i$ and $\bar q=\abs{\bar\sigma}^2/N$, we have $\bar q$ close to $Q+(1-Q)/n$, since the diagonal of the overlap matrix enters the mean. The point $u=\bar\sigma/\sqrt{\bar q}$ lies on $S_N$, and $H_Q(u)=H_N(\sqrt{Q/\bar q}\,\bar\sigma)$. We therefore choose $n$ large, so that $\bar q$ is close to $Q$, but fixed before $N\to\infty$. A localization lemma then moves $u$ to a nearby critical point of $H_Q$, and the given sample lies in its band (for pure models, in its cap of a fixed latitude). In the pure proof, compactness gives a number of additional replicas that is uniform over the limiting overlap laws. In the mixed proof, this number is chosen after passing to a limiting overlap law and may depend on that subsequence; the subsequence argument then proves the full-sequence coarse-cover limit. In both cases, it is fixed before the dimension tends to infinity.

For non-pure mixtures, the coarse band family can contain exponentially many critical points, and the next two inputs pass to a fixed window. Let $Z_H(D)=\int_De^{\beta H}\dd\nu_N$. Suppose that a center $u$ in the window $[m_N-L,m_N+L]$ satisfies $Z_H(B_u)\ge\delta D_Ne^{\beta(H_Q(u)-m_N)}$, for a deterministic scale $D_N$. Then $Z_{N,\beta}\ge\delta e^{-\beta L}D_N$, and on the event that $\max H_Q\le m_N+M$, we have
\begin{equation}\label{eq:ov-capture-split}
\pi\bigl(B_{N,M}^c\bigr)\le\pi\Bigl(\Bigl[\bigcup_{u\text{ coarse}}B_u\Bigr]^c\Bigr)+\frac{e^{\beta L}}{\delta D_N}\sum_{\substack{u\text{ coarse}\\H_Q(u)<m_N-M}}Z_H(B_u)
\end{equation}
(\cref{rem:mb-assembly}). In the second step, we bound the expected sum. In the third step, we show that light bands, those with $Z_H(B_u)<\delta D_Ne^{\beta(H_Q(u)-m_N)}$, are rare, so that a center as above exists.

For the bands of centers below the window, we use the Kac--Rice formula with the band partition function as a weight. We show that the Palm expectation of $Z_H(B_u)$ at a center of value $E$, that is, its conditional mean weighted by the determinant, is at most $C_1D_Ne^{\beta(E-m_N)+C_2N(E/N-E_0)^2}$ (\cref{prop:mb-upper}). The density of centers at value $m_N-v$ is at most $Ce^{c_+v}$, with $c_0<c_+<\beta$ (\cref{lem:mb-envelope}). We choose the coarse energy width $0<\tau\le\tau_0=\min\{\varepsilon_0,(\beta-c_+)/(4C_2)\}$, as in \cref{cor:mb-low-sum}. This absorbs the quadratic correction into the decay in the energy deficit, and the expected sum in \eqref{eq:ov-capture-split} is at most $CD_Ne^{-a_0M}$, with $a_0=(\beta-c_+)/2>0$. This is the rigorous form of the heuristic $e^{c_0v}e^{-\beta v}$ above, and it is the only place where $c_0<\beta$ enters the concentration argument for non-pure mixtures. The difficulty is that the upper bound must hold without any power of $N$, since $M$ is fixed while $N\to\infty$. We write the Palm expectation as the ratio of a Laplace integral over the radial derivative and the latitude to a Laplace integral over the radial derivative alone. In the ratio, one latitude width $N^{-1/2}$ remains, and it cancels the growth of the constant $(N/2\pi)^{1/2}$ in the coarea formula over latitudes. The determinant prefactors are bounded (\cref{lem:mb-prefactor,lem:mb-laplace}).

The bound on light bands rests on the structure of the field on the band. At latitude exactly $\sqrt Q$, the field is the prescribed value plus a residual field with covariance $N\psi_Q$, whose law does not depend on the marks. Since the residual model is strictly replica symmetric with $\beta^2\psi_Q''(0)<1$, the theorem of Huang, Montanari, and Pham \cite[Lemma~3.2 and Remark~3.3]{HuangMontanariPham2024} implies that its normalized partition function is tight away from zero (\cref{lem:pre-HMP-Z}). We pass from this latitude to the band with Jensen's inequality over a symmetric latitude interval of width $N^{-1/2}$ (\cref{lem:mb-jensen}). The Hessian determinant is correlated with the field on the band. We transfer the bound to the Palm law with the normalized second moment of the determinant, which needs no information on this correlation (\cref{prop:mb-lower}). Then the expected number of light bands in each fixed window is small (\cref{cor:mb-bad}), which is the remaining hypothesis of \cref{prop:red-capture}.

For pure models, the radial derivative is a multiple of the value, $K_Q=pH_Q$, and the conditioning involves the value and the gradient only. We show that the band partition functions at a critical point have normalized second moments bounded by a constant and a lower tail of order $1/T$ at level $e^{-T}$. The proof is an interpolation argument that needs an exponential moment of the covariance of the field on the band (\cref{lem:pc-wick,prop:pc-bands}). This moment is finite because the residual coefficient satisfies $\beta^2\psi_Q(1)<1/2$ for every $p\ge4$ and $\beta>\bs(p)$ (\cref{lem:th-pure-residual}). The contribution of the bands around critical points with a large energy deficit decays exponentially in the deficit, at a rate that is positive because $x<1$ and the cap latitude is chosen close to $\sqrt Q$ (\cref{lem:pc-saddle} and the proof of \cref{thm:pc-capture}).

\subsection{The conditioned hitting estimate}\label{sec:ov-entrance}

The hitting estimate for the caps under critical-point conditioning, \cref{thm:ent-main}, is the main dynamical input of \cref{part:condensed}. For a finite mixture of minimum degree $d\ge4$, fixed $Q\in(0,1]$, $\beta\ge0$, and fixed bounds $\bar h,\bar k$ on the marks, there exists $a_*>0$ such that for every $a\in(0,a_*]$ there are $\gamma,c,C>0$ for which, for all sufficiently large $N$,
\begin{equation}\label{eq:ov-entrance}
\E_{h,k,u}\Bigl[\Pp^{H_N}_{\nu_N}\Bigl(\sup_{0\le t\le e^{\gamma N}}\abs{R(X_t,u)}\ge a\Bigr)\Bigr]\le C\bigl(N^{1-d/2}+N^{3-d}\bigr)+e^{-cN},
\end{equation}
uniformly in $u$ and admissible marks $(h,k)$, meaning $\abs h\le\bar h$, $\abs k\le\bar k$, and $k=ph$ when the mixture is pure of degree $p$. Here the diffusion is the original Langevin diffusion, run in the conditioned field. \Cref{sec:ent-mechanism} explains the proof in detail; here we describe its main points.

We begin with the unconditioned field. For an axis $u$ that is uniform on $S_N$ and independent of the field, the mean Gibbs mass of $C_a(u)$ is $\int\nu_N\{u:\abs{R(x,u)}\ge a\}\,\pi(\dd x)$. The rotation invariance of $\nu_N$ implies that this equals $\nu_N(C_a(v))$ for any fixed $v\in S_N$, which is exponentially small, for every realization of the field. The conditioning changes the whole field. By Gaussian regression, the conditioned field is $H_N+\delta_{h,k,u}$, where the correction is explicit and of order $N$ near $u$ (\cref{prop:ent-regression}). We cannot treat it as a small perturbation. However, it vanishes to high order at the equator of $u$. The following correction bound is on the regularity event $\Omega_N$ of \cref{lem:ent-regularity}, whose complement has probability at most $e^{-N}$: with $r=R(x,u)$ and $V_x(u)=N^{-1/2}\nabla_{S_N}H_Q(u)\cdot x$, we have
\begin{equation}\label{eq:ov-correction}
\abs{\delta_{h,k,u}(x)}\le CN\abs r^d+C\sqrt N\abs r^{d-1}\abs{V_x(u)}.
\end{equation}
A fixed point $x$ has overlap of order $N^{-1/2}$ with a uniform axis, and $V_x(u)$ is of order one. We see that the correction is of order $N^{1-d/2}$, and this is the source of both polynomial errors in \eqref{eq:ov-entrance}. Fix the field in $\Omega_N$. For a sufficiently small fixed $\alpha>0$, let $\chi$ be an even smooth cutoff equal to one on $[-1,1]$ and zero outside $[-2,2]$, set $s=\nabla_{S_N}H_Q(u)\cdot x/N$, and define $\delta_\alpha(u,x)=\chi(r/\alpha)\chi(s/\alpha)\delta_{h,k,u}(x)$ and $w(u,x)=\beta\delta_\alpha(u,x)$. We compare the dynamics in this truncated field, at a uniform axis, with the unconditioned dynamics. Only two averaged moments enter, and they hold uniformly in $x$, the field in $\Omega_N$, and the admissible marks (\cref{prop:ent-multipliers}):
\begin{equation}\label{eq:ov-moments}
\eta=\E_u\abs{e^w-1}=O\bigl(N^{1-d/2}\bigr),\qquad\epsilon=\E_u\bigl[e^w(1-e^{-w})^2\bigr]=O\bigl(N^{2-d}\bigr).
\end{equation}
There is an apparent difficulty here, since the point $x$ in these averages may depend on the field, and so does the gradient $\nabla_{S_N}H_Q(u)$ in $V_x(u)$. The way out is to fix the field and $x$ and to average $u$ over each latitude $\{u:R(x,u)=r\}$. On a latitude, $V_x$ concentrates by the logarithmic Sobolev inequality on the sphere \cite{BakryEmery1985}, and the Funk--Hecke formula \cite{AtkinsonHan2012} gives its mean exactly (\cref{lem:ent-lsi,lem:ent-latitude,prop:ent-multipliers}).

The main obstacle is the length $e^{\gamma N}$ of the time interval, since a comparison of path laws, for instance by Girsanov's formula, costs an amount proportional to the time. To overcome this obstacle, we compare resolvents instead. Let $\tau$ be the killing time for a killing rate $\kappa$ supported near the caps. The entropy--resolvent comparison of \cref{lem:ent-resolvent} bounds the Laplace transform $\E e^{-\lambda\tau}$ for the perturbed process by the integrated killing rate against a resolvent of the unperturbed process, plus a cost
\begin{equation}\label{eq:ov-cost}
\eta+\sqrt\epsilon+\epsilon\,(1+\mathcal H_0)=O\bigl(N^{1-d/2}+N^{3-d}\bigr),
\end{equation}
where $\mathcal H_0$ is an entropy of the density of $\nu_N$ with respect to the Gibbs measure. Note that $\mathcal H_0$ is of order $N$, since $\abs{H_N}\le KN$ with overwhelming probability. The point is that the cost is not multiplied by the time: with $\lambda=e^{-\gamma N}$, the factor $e^{\gamma N}$ multiplies only the integrated killing rate, which is exponentially small when $\gamma$ is small. The comparison cost other than the killing term depends on the initial density $f_0$ through $\mathcal H_0$; the killing term also involves its unperturbed resolvent $\lambda(\lambda+A_0)^{-1}f_0$, where $A_0$ is the generator of the unperturbed dynamics. A commutator identity in the proof cancels a cross term that would otherwise be of order one. For $d=4$, both terms in \eqref{eq:ov-cost} are of order $N^{-1}$. For $d=3$, the second term is of order one, and this is where we need the hypothesis on the minimum degree.

Two further steps return to the true process. Recall that we have truncated the correction at a fixed scale $\alpha$ in both $r$ and $s$, since the comparison requires a bounded perturbation. The truncated correction is at most $CN\alpha^d$, while the cap of latitude $\alpha/4$ has mass $e^{-cN\alpha^2}$. Since $d>2$, we can choose $\alpha$ small, and then $\gamma$ small. Except with exponentially small probability, a hit of the cap forces the smooth killing rate to act during a time interval of fixed length, and this converts the killing time into a hitting time (\cref{prop:ent-truncated}). The true and truncated processes coincide until $\abs r$ or $\abs s$ reaches $\alpha$. Reaching $\abs r=\alpha$ requires crossing the cap of latitude $a=\alpha/2$ first. Reaching $\abs s=\alpha$ is exponentially unlikely, because the gradient $\nabla_{S_N}H_Q(u)$ is exactly the Gaussian vector that the conditioning removes by regression. In the unrestricted regression coupling, conditionally on $u$, this tangent gradient is independent of the conditioned field and its whole trajectory. Conditionally on $u$ and the trajectory, $s_t$ is centered Gaussian with variance at most $a_Q/N$ at each fixed time. The resulting exponential bound for the gradient cutoff also holds after restricting to $\Omega_N$, by bounding the unrestricted probability and dividing by $\Pp(\Omega_N)$ (\cref{lem:ent-boundary}). Finally, we pass from a uniform axis to every deterministic axis by rotation invariance.

\subsection{Other initial laws}\label{sec:ov-initial}

We deduce \cref{thm:main-initial} from \cref{thm:main-condensed} and the polynomial bound of \cref{prop:cp-transfer}, by three arguments that use only the rotation invariance of the disorder (\cref{sec:red-initial}). For part (i), let $\mu$ be a deterministic initial law. The expected distance $\mu\mapsto\E\norm{\mu P_T-\pi}_{\TV}$ is convex, and it is invariant under rotations of $\mu$, because we can compensate a rotation of the initial law by a rotation of the field, whose law is invariant. The rotation average of $\mu$ is $\nu_N$, and convexity implies that
\begin{equation}\label{eq:ov-haar}
\E\bigl[1-\norm{\mu P_T-\pi}_{\TV}\bigr]\le\E\bigl[1-\norm{\nu_NP_T-\pi}_{\TV}\bigr]
\end{equation}
(\cref{thm:red-haar}). This is an instance of the classical symmetrization over a compact group, the same averaging that underlies the Hunt--Stein theorem \cite[Section~8.5]{LehmannRomano2005}. In particular, $\nu_N$ is the hardest case among the initial laws chosen independently of the disorder.

For parts (ii) and (iii), the same invariance shows that the expected mass of the band union at time $T_N$ is the same from every starting point (\cref{lem:red-equal-mass}). \Cref{prop:cp-transfer} implies that this mass is at most $q_{N,M}\le C_{M,r}N^{-(d-2)(1-1/r)/2}$. We handle a field-dependent choice among $K_N\le C_0N^\kappa$ prescribed laws by a union bound, since $K_Nq_{N,M}\to0$ when $\kappa<(d-2)/2$ and $r$ is large (\cref{cor:red-selection}). The loss of the factor $1-1/r$ in the H\"older step is why we exclude the endpoint $\kappa=(d-2)/2$. We handle a field-dependent law $\mu$ with small relative entropy with respect to $\nu_N$ by binary data processing. Let $q$ be the mass of the band union under $\nu_NP_{T_N}$. Then its mass under $\mu P_{T_N}$ is at most $(\KL(\mu\Vert\nu_N)+\log2)/\log(1/q)$ (\cref{lem:pre-binary,cor:red-entropy}). For any fixed $r>1$ and $0<\tau<(d-2)(1-1/r)/2$, we have $q\le N^{-\tau}$ with probability tending to one, by Markov's inequality. The entropy bound then tends to zero in disorder probability when $\KL(\mu\Vert\nu_N)/\log N\to0$ in that probability mode.

\section{Below the static threshold}\label{sec:ov-prestatic}

Let $\bd<\beta<\bs$, and recall the plateau value $q_*$ of \eqref{eq:intro-qstar}. The proof of \cref{thm:main-memory} has a probabilistic part, in which we describe the Gibbs measure as seen from a Gibbs-typical point, and a deterministic part, in which we build one observable from this description and derive the five conclusions from it.

\subsection{The constrained pressure and the equality at \texorpdfstring{$q_*$}{q*}}\label{sec:ov-pressure}

For a reference point $x$ and an overlap $r$, consider the Gibbs mass of the configurations $u$ with $R(x,u)$ close to $r$. For a reference point $x$ drawn from $\pi$, the constrained pressure $F_{\rm FP}$ of \cref{def:pr-FP} is the Franz--Parisi potential written as a pressure. With $F=\beta^2\xi(1)/2$, for a fixed $h>0$ such that $[r-h,r+h]\subset(-1,1)$, the leading exponential form under the Gibbs-reference law is
\begin{equation}\label{eq:ov-FP}
\pi\bigl\{u:R(x,u)\in[r-h,r+h]\bigr\}\approx\exp\Bigl\{N\Bigl[\max_{v\in[r-h,r+h]}F_{\rm FP}(v)-F\Bigr]\Bigr\}.
\end{equation}
The midpoint exponent $F_{\rm FP}(r)-F$ is recovered by sending $h$ to zero after the limit in $N$ (\cref{prop:pr-cross}).
A direct computation is difficult, because $x$ depends on the disorder. Following El Alaoui, Montanari, and Sellke \cite{ElAlaouiMontanariSellke2025}, who treated pure models, we work under a planted law, in which $x$ is uniform on $S_N$ and the field is a Gaussian field plus the drift $\beta N\xi(R(x,\cdot))$. Below the static threshold, events of exponentially small probability under the planted law have exponentially small probability under the joint law of the field and a reference point drawn from $\pi$ (\cref{lem:pr-transfer}). Let $\eta_r(t)=\xi(r^2+(1-r^2)t)-\xi(r^2)$. Under the planted law, with $n=N-1$, the field on the section $\{u:R(x,u)=r\}$ has the law of $\beta N\xi(r)+W+Y$, where $W$ is centered Gaussian with variance $N\xi(r^2)$ and is independent of the transverse spherical field $Y$, whose covariance is $N\eta_r(R_n)=n(N/n)\eta_r(R_n)$ (\cref{lem:pr-section}). The Parisi formula then gives
\begin{equation}\label{eq:ov-cross}
F_{\rm FP}(r)=\beta^2\xi(r)+\frac12\log(1-r^2)+\cP_\beta(\eta_r)
\end{equation}
(\cref{prop:pr-cross}). For $r\ne0$, the mixture $\eta_r$ has a positive term of degree one, even though $\xi$ has none; at $r=0$, we have $\eta_0=\xi$. The Parisi formula of Chen \cite{Chen2013} allows the degree-one term.

The infimum $\cP_\beta(\eta_r)$ is not explicit, but an upper bound suffices. For the trial bound, we restrict to $r\in[0,1)$. Evaluating the Crisanti--Sommers functional of $\eta_r$ at the point mass $\delta_{r/(1+r)}$, we obtain $F_{\rm FP}(r)\le\Psi(r)$ (\cref{lem:pr-trial}), where
\begin{equation}\label{eq:ov-trial}
\Psi(r)=F+\frac12\Phi_\beta(r),\qquad\Psi'(r)=\frac r{2(1-r)}\bigl(\beta^2A(r)-1\bigr).
\end{equation}
For pure models, this is the trial bound of \cite[Proposition~3.9]{ElAlaouiMontanariSellke2025}. By the definition of $q_*$, the trial pressure increases on an interval $(q_*-\eps_0,q_*)$ and does not increase on $[q_*,1)$. The new point is that $F_{\rm FP}(q_*)=\Psi(q_*)$ for every finite mixture (\cref{prop:pr-certificate}). This amounts to showing that $\delta_{r_0}$, with $r_0=q_*/(1+q_*)$, minimizes the Crisanti--Sommers functional of $\eta_{q_*}$, which we check with the Crisanti--Sommers optimality criterion. Above $r_0$, the derivative of the first variation has the sign of $\beta^2A-1$ to the right of $q_*$, which is nonpositive. Below $r_0$, we need that $\xi'(v)/(v-q_*^2)$ is nonincreasing on $(q_*^2,q_*]$. We deduce this from the monotonicity of the ratio $\chi(v)=v\xi''(v)/\xi'(v)$, whose derivative is the variance of $p-1$ under the weights $pc_pv^{p-1}/\xi'(v)$, divided by $v$, together with $A'(q_*)\le0$. We allow $q_*$ to be a degenerate root of the plateau equation.

Fix an admissible collar $0<q_1<q_2<q_*$, so that $\beta^2A(r)>1$ on $(q_1,q_*)$. Since $\Psi$ increases on $[q_1,q_*]$ and does not increase on $[q_*,1)$, the equality at $q_*$ implies two bounds. On the cap $[q_1,1)$, the constrained pressure is at most $F-J$, with equality at $q_*$, where $J=-\Phi_\beta(q_*)/2>0$. On the collar $[q_1,q_2]$, it is at most $F-J-G$, where $G=\Psi(q_*)-\Psi(q_2)>0$. After fixing the collar, choose $0<\delta<\min\{J,G\}/8$. Using the planted law and likelihood transfer, we turn these limits into quenched estimates (\cref{prop:pr-good}). With probability at least $1-Ce^{-aN}$, there exists a set $\cG$ of good reference points with $\pi(\cG^c)\le e^{-aN}$ such that, for every $x\in\cG$,
\begin{equation}\label{eq:ov-good}
\pi\{u:R(x,u)\ge q_2\}\ge e^{-N(J+2\delta)},\qquad\pi\{u:q_1\le R(x,u)\le q_2\}\le e^{-N(J+G-2\delta)}.
\end{equation}
In words, around a Gibbs-typical point there is a region of Gibbs mass about $e^{-NJ}$ at overlap close to $q_*$, and it is separated from the rest of the Gibbs measure by a collar of exponentially smaller mass.

The second input is a uniform bound on the Gibbs masses of caps. By the replica concentration of Huang, Montanari, and Pham \cite[Lemma~7.7]{HuangMontanariPham2024}, for each fixed $s>0$, two independent samples from $\pi$ have overlap larger than $s$ in absolute value with Gibbs probability at most $2e^{-cN}$, on an event of probability at least $1-e^{-cN}$ (\cref{prop:pr-replica}). We convert this into a bound for every cap, uniformly over its axis, with an elementary barycenter estimate (\cref{lem:pr-cap}):
\begin{equation}\label{eq:ov-cap}
\sup_{v\in S_N}\pi\{x:R(x,v)\ge s\}\le\Bigl[\frac{2-s^2}{s^2}\,(\pi\otimes\pi)\bigl\{R>s^2/2\bigr\}\Bigr]^{1/2}.
\end{equation}
The axis may depend on the field, and we need no union bound over axes. In the same way, we find that $\abs{\pi(U)}^2\le s+e^{-cN}$ (\cref{cor:pr-cap}).

\subsection{The overlap-weighted Gibbs mean}\label{sec:ov-barycenter}

The estimates \eqref{eq:ov-good} suggest an observable that records the region around $x$. Let $h\colon[-1,1]\to[0,K]$ be smooth and nonincreasing, with $h=K$ on $[-1,q_1]$ and $h=0$ on $[q_2,1]$, and let $\psi_u(x)=e^{-Nh(R(x,u))}$. The overlap-weighted Gibbs mean (\cref{def:sp-barycenter}) is
\begin{equation}\label{eq:ov-barycenter}
B_N(x)=\int U(u)\,w_x(\dd u),\qquad w_x(\dd u)=\frac{\psi_u(x)^2\,\pi(\dd u)}{\int\psi_{u'}(x)^2\,\pi(\dd u')}.
\end{equation}
In words, $B_N(x)$ is the mean of $U$ under the Gibbs measure reweighted by a cutoff that suppresses the points $u$ with $R(x,u)<q_1$. We call $w_x$ the overlap-reweighted Gibbs measure seen from $x$. Note that $\abs{B_N}\le1$.

We show that $B_N$ has two properties at the points $x\in\cG$, and both follow from \eqref{eq:ov-good} (\cref{lem:sp-barycenter}). The first is that its gradient is small. Differentiating \eqref{eq:ov-barycenter}, we find that $\nabla B_N(x)$ is twice the $w_x$-mean of $(U(u)-B_N(x))\otimes\nabla\log\psi_u(x)$. The factor $\nabla\log\psi_u(x)$ vanishes unless $u$ lies in the collar of $x$, and only the $w_x$-mass of the collar counts. With $d_N(x)$ the denominator in \eqref{eq:ov-barycenter}, we obtain
\begin{equation}\label{eq:ov-gradient}
\norm{\nabla B_N(x)}_{\rm HS}^2\le4N\norm{h'}_\infty^2\,\frac{\pi\{u:q_1\le R(x,u)\le q_2\}}{d_N(x)}\le4N\norm{h'}_\infty^2e^{-N(G-4\delta)},
\end{equation}
where we used $d_N(x)\ge\pi\{u:R(x,u)\ge q_2\}$ and \eqref{eq:ov-good}. The second property is alignment. The references with $R(x,u)\le q_1$ carry $w_x$-mass at most $e^{-2KN}/d_N(x)$, and with $2K=J+2\delta+1$ we find that $U(x)\cdot B_N(x)\ge q_1-(1+q_1)e^{-N}$ (\cref{sec:mem-proof}).

The definition averages over all references $u$, including those outside $\cG$, at which the estimates \eqref{eq:ov-good} may fail. The function $B_N$ is then smooth on all of $S_N$, and $\norm{\nabla B_N}_{\rm HS}^2\le4N\norm{h'}_\infty^2$ at every point, because $\abs{\nabla\psi_u}^2\le N\norm{h'}_\infty^2\psi_u^2$ and $w_x$ is a probability measure (\cref{lem:sp-barycenter}). The points $x$ outside $\cG$ therefore enter the vector energy only through their Gibbs mass:
\begin{equation}\label{eq:ov-energy}
\cE_{\rm vec}(B_N)\le4N\norm{h'}_\infty^2\bigl(e^{-N(G-4\delta)}+e^{-aN}\bigr).
\end{equation}

\subsection{From one observable to five conclusions}\label{sec:ov-five}

On the event of \cref{prop:pr-good,cor:pr-cap}, the rest of the proof is deterministic (\cref{sec:spectral,sec:memory}). We list in \cref{tab:ov-five} the tool that we use for each conclusion of \cref{thm:main-memory}.

\begin{table}[ht]
\centering
\small
\begin{tabular}{@{}>{\raggedright\arraybackslash}p{0.27\textwidth}>{\raggedright\arraybackslash}p{0.4\textwidth}>{\raggedright\arraybackslash}p{0.25\textwidth}@{}}
\toprule
Conclusion & Tool & Result\\
\midrule
(i) many slow modes & averaged variational principle applied to translated bumps of $B_N$; bound on caps & \cref{lem:sp-frame,lem:sp-bumps}\\
(ii) projection onto slow modes & spectral projection of $Z_N$ against $B_N-\pi(B_N)$; small magnetization & \cref{lem:sp-projection}\\
(iii) correlation plateau & spectral calculus, from (ii) & \cref{lem:sp-projection}\\
(iv) Gibbs-typical starts & displacement bound and Markov's inequality; bound on caps & \cref{lem:sp-trap}\\
(v) endpoint information & binary data processing and the chain identity, from (iv) & \cref{lem:pre-binary,lem:pre-MI}\\
\bottomrule
\end{tabular}
\caption{The five conclusions of \cref{thm:main-memory} and the deterministic results that derive them from the energy and the alignment of the overlap-weighted Gibbs mean.}
\label{tab:ov-five}
\end{table}

For (i), we use the averaged variational principle of Harrell and Stubbe \cite[Theorem~3.1]{HarrellStubbe2014} in the following form (\cref{lem:sp-frame}). Suppose that a family of nonnegative trial functions $f_z$, indexed by a measure space, has total squared $L^2(\pi)$ mass $0<\mathsf T=\int\norm{f_z}_2^2\,\mu(\dd z)<\infty$ and averaged energy at most $r\mathsf T$, that each function satisfies $\pi(f_z)\le\mathsf R$, and that the family covers each point with total weight at most $\mathsf C$. Then the number of eigenvalues below $\Lambda>r$ is at least $(1-r/\Lambda)\mathsf T/(\mathsf R\mathsf C)$. We apply this to the bumps $f_z(x)=\alpha(N[1-\abs{B_N(x)-z}^2/\rho^2])$, with $z$ averaged over a ball of $\R^N$ (\cref{lem:sp-bumps}). A bump is supported where $B_N$ is within $\rho$ of $z$. Since $B_N$ is aligned with $U$ on $\cG$, the part of this set in $\cG$ lies in a cap of a fixed latitude, and \eqref{eq:ov-cap} bounds its Gibbs mass by $e^{-cN}$. With the index measure of \cref{lem:sp-bumps}, we have $\int f_z(x)\,\mu(\dd z)\le1$ at every $x$ and $1/2\le\int\norm{f_z}_2^2\,\mu(\dd z)\le1$. The averaged energy of the bumps is at most $CN^2\rho^{-2}\cE_{\rm vec}(B_N)$, which is exponentially small by \eqref{eq:ov-energy}. This shows that at least $e^{cN}$ eigenvalues lie below $e^{-2\gamma N}$. Unlike the smoothed indicators of disjoint clusters used in \cite{ElAlaouiMontanariSellke2025}, the bumps overlap, and the averaged principle needs only the bounds $\mathsf R$ and $\mathsf C$.

For (ii) and (iii), let $Z_N=U-\pi(U)$, $\overline B_N=B_N-\pi(B_N)$, and $\Pi_\Lambda=\ind_{(0,\Lambda]}(-L)$. The part of $\overline B_N$ above $\Lambda$ has squared norm at most $\cE_{\rm vec}(B_N)/\Lambda$, and we find that
\begin{equation}\label{eq:ov-projection}
\norm{\Pi_\Lambda Z_N}_2\,\norm{\overline B_N}_2\ge\ip{Z_N}{\overline B_N}-\bigl(\cE_{\rm vec}(B_N)/\Lambda\bigr)^{1/2}
\end{equation}
(\cref{lem:sp-projection}). The alignment implies that $\ip U{B_N}\ge q_1-o(1)$, and the magnetization $\pi(U)$ is small by \cref{cor:pr-cap}. Together, they give $\ip{Z_N}{\overline B_N}\ge q_0$ for a fixed $q_0<q_1$ close to $q_1$, while $\norm{\overline B_N}_2\le1$. With $\Lambda=e^{-2\gamma N}$, we obtain $\norm{\Pi_\Lambda Z_N}_2^2\ge q_0^2-o(1)$. For every $b<q_*^2$, we can choose the collar and $q_0$ so that $q_0^2>b$, and this proves (ii). The connected correlation satisfies $\ip{Z_N}{P_tZ_N}\ge e^{-\Lambda t}\norm{\Pi_\Lambda Z_N}_2^2$, which implies (iii) up to time $e^{\gamma N}$.

Conclusion (iv) follows from \eqref{eq:ov-endpoint-energy} and Markov's inequality. Let $T=e^{\gamma N}$. For all $x$ outside a set of Gibbs mass $e^{-cN}$, we find that the law $P_T(x,\cdot)$ gives mass at least $1-e^{-\gamma N}$ to the set $D_x=\{y:\abs{B_N(y)-B_N(x)}\le\rho\}$. For $x\in\cG$, the set $D_x$ lies in the union of $\cG^c$ and a cap around the direction of $B_N(x)$, and its Gibbs mass is at most $w=e^{-cN}$. Testing $D_x$ at time $T$, we obtain $\norm{P_T(x,\cdot)-\pi}_{\TV}\ge1-e^{-\gamma N}-w$, and the contraction of total variation gives the same lower bound for every $0\le t\le T$. For (v), by binary data processing at time $T$, we have $\KL(P_T(x,\cdot)\Vert\pi)\ge(1-e^{-\gamma N})\log(1/w)-\log2$. The contraction of relative entropy gives this bound for every $0<t\le T$, and it is at least $cN$ for all sufficiently large $N$ (\cref{lem:sp-trap}). We conclude with the chain identity $I(X_0;X_t)=\int\KL(P_t(x,\cdot)\Vert\pi)\,\pi(\dd x)$ of \cref{lem:pre-MI}, which goes back to the golden formula \cite[Theorem~4.1]{PolyanskiyWu2025}.

\subsection{Correlation bounds at every level below \texorpdfstring{$q_*$}{q*}}\label{sec:ov-plateau}

The proof above gives levels below $q_*^2$. The physics prediction uses the largest positive solution of the plateau equation, which equals $q_*$ except when $\beta^2A$ touches one from below at an isolated point farther to the right (see \cref{sec:intro-model}). \Cref{thm:main-plateau} gives lower bounds approaching $q_*$. The loss of the square occurs in \eqref{eq:ov-projection}, where we used $\norm{\overline B_N}_2\le1$. Recall from \eqref{eq:ov-barycenter} that $\abs{B_N(x)}^2$ is the mean overlap of two independent samples from $w_x$, which concentrate at overlap about $q_*$ with $x$. If their mutual overlap also concentrates at $q_*$, then $\norm{\overline B_N}_2^2\le q_*+o(1)$. Combining this with \eqref{eq:ov-projection}, for a fixed sufficiently small $0<\vartheta<\min\{1,q_1/3\}$ and a time exponent $\gamma$ chosen as in \cref{sec:pl-reduction}, the norm, magnetization, and energy estimates give, with probability tending to one,
\begin{equation}\label{eq:ov-sharper}
\norm{\Pi_{\Lambda_N}Z_N}_2^2\ge\frac{(q_1-3\vartheta)^2}{q_*+\vartheta},\qquad\Lambda_N=e^{-2\gamma N}.
\end{equation}
Choosing $q_1$ close to $q_*$ and then $\vartheta$ small makes this lower bound exceed any fixed $b<q_*$. The main step of \cref{sec:plateau} is therefore the estimate $\int\abs{B_N}^2\dd\pi\to q_*$ in probability (\cref{prop:pl-norm}).

We prove it under the planted law. We write each $u$ through its latitude $r=R(x,u)$ and its transverse part $y$, a point of the sphere in $x^\perp$. For two such points, we have $R(u,u')=rr'+((1-r^2)(1-r'^2))^{1/2}R_n(y,y')$. The latitude concentrates at $q_*$ under $w_x$, by the constrained pressure. It remains to show that the transverse overlap $R_n(y,y')$ concentrates at $t_0=q_*/(1+q_*)$, since $q_*^2+(1-q_*^2)t_0=q_*$. By the optimality of the replica-symmetric trial at $q_*$, the Parisi measure of the transverse spin glass at latitude $q_*$ is the point mass at $t_0$ (\cref{lem:pl-profile}). Note that the transverse covariance $\eta_{q_*}$ contains the monomials $t$ and $t^2$ with positive coefficients. We multiply two Gaussian components of the field that contribute to these monomials by amplitudes $s_1$ and $s_2$ and differentiate the expected logarithm of the reweighted partition function in these amplitudes. Using Gaussian integration by parts, we find that the derivative in $s_k$ at $s=(1,1)$ is, once the latitude has concentrated, an affine function of the $k$-th moment of the transverse overlap (\cref{lem:pl-ibp}). We compute its limit from the Crisanti--Sommers formula and the envelope theorem \cite{MilgromSegal2002}: it equals $\beta^2a_k^2(1-t_0^k)$, where $a_k^2>0$ is the coefficient of $t^k$ contributed by the $k$-th component (\cref{lem:pl-envelope}). Convexity in the amplitudes lets us pass from the convergence of the pressures to the convergence of their derivatives \cite[Theorem~24.5]{Rockafellar1970}. We find that the first two moments of the transverse overlap tend to $t_0$ and $t_0^2$, and that its variance tends to zero. This use of derivatives in the coefficients of a covariance is classical in the theory of mean-field spin glasses \cite{Panchenko2008}; see \cite[Proposition~2.2]{AuffingerElAlaouiSellke2025} for spherical mixtures. Finally, since $(\xi,\beta)$ is strictly replica symmetric with $\beta^2\xi''(0)<1$, the planted law and the Gibbs-reference law are mutually contiguous \cite[Lemma~3.2, Remark~3.3, and Corollary~3.5]{HuangMontanariPham2024}, which transfers the estimate to the original disorder. Neither the convex derivative convergence nor the likelihood transfer supplies a rate. The derivative argument does not give exponentially small planted failure probabilities, so the exponential transfer of \cref{lem:pr-transfer} cannot be applied; this is why (ii) and (iii) of \cref{thm:main-plateau} are proved only with probability tending to one.

\subsection{Eigenvalues near the dynamical threshold}\label{sec:ov-cusp}

The proof of \cref{thm:main-cusp} reuses the estimates above.

For \cref{thm:main-cusp}, we count slow modes with smooth cap cutoffs $f_u(x)=\varphi(R(x,u))$ as trial functions, where $\varphi$ is smooth, vanishes on $(-\infty,q_1]$, and equals one on $[q_2,\infty)$, indexed by good references with measure $\ind_{\cG}(u)\pi(\dd u)$ (\cref{rem:mem-direct-frame,lem:inf-frame}). The exponent of the eigenvalue bound can then be any number below the collar gap $G$, which we compute as $\beta\downarrow\bd$. Near the largest maximizer $q_d$ of $A$, we write $\eps=\beta^2A_{\max}-1$ and $r=q_d+hz$ with $h=(\eps/a)^{1/(2m)}$. On the component of $\{\beta^2A>1\}$ that ends at $q_*$, \eqref{eq:ov-trial} shows that, heuristically, $\Psi'(r)\approx q_d\eps(1-z^{2m})/(2(1-q_d))$. Integrating over $z\in[-1,1]$, we find
\begin{equation}\label{eq:ov-cusp}
\Psi(q_*)-\Psi(q_-)\approx\frac{q_d}{2(1-q_d)}\,h\eps\int_{-1}^1(1-z^{2m})\dd z=C_d\,\eps^{1+1/(2m)},
\end{equation}
where $q_-$ is the left end of this component. The precise forms are \eqref{eq:inf-barrier} and, for the collar gap, \eqref{eq:inf-collar-gap}. In the counting bound, the core exponent $J$ enters both $\mathsf T$ and $\mathsf R$ and cancels, and the covering constant $\mathsf C$ is the largest Gibbs mass of a cap. We use a cap bound whose exponent is uniform in the temperature (\cref{lem:inf-replica}). The number of slow modes then stays $e^{\kappa_*N}$, with $\kappa_*$ independent of $\eps$, while the exponent of the eigenvalue bound tends to zero. Here $\eps>0$ is fixed before $N\to\infty$; the probability exponent and the threshold for $N$ may depend on $\eps$ and $\vartheta$. The computed barrier is that of the trial pressure, and it yields a lower bound for the constrained-pressure gap, rather than an asymptotic formula for that gap or for the eigenvalues.

\section{Dependencies among the main steps}\label{sec:ov-dependencies}

We list in \cref{tab:ov-condensed,tab:ov-prestatic} the principal proof inputs for each main result, including selected intermediate results, in the order discussed above. We suppress most general preliminary facts, but retain the Crisanti--Sommers optimality criterion where it is a central variational input.

\begin{table}[ht]
\centering
\footnotesize
\begin{tabular}{@{}>{\raggedright\arraybackslash}p{0.21\textwidth}>{\raggedright\arraybackslash}p{0.30\textwidth}>{\raggedright\arraybackslash}p{0.41\textwidth}@{}}
\toprule
Result & Content & Uses\\
\midrule
\cref{thm:main-condensed} & separation from the uniform law & \cref{prop:cp-transfer}, \cref{prop:th-structure,lem:th-superstatic}; counts and concentration \cref{prop:pc-counts,thm:pc-capture} in pure models, or \cref{thm:mb-capture} with the counts of \cref{thm:mc-counts,lem:mb-counts} in non-pure mixtures\\
\cref{thm:main-initial} & other initial laws & \cref{thm:main-condensed}, \cref{prop:cp-transfer}, \cref{thm:red-haar}, \cref{cor:red-selection}, \cref{cor:red-entropy}\\
\cref{prop:cp-transfer} & a common time horizon for all fixed windows & \cref{prop:red-consumer}, \cref{cor:ent-rectangle}, \cref{thm:kr-transfer}\\
\cref{thm:kr-transfer} & determinant-weighted transfer & \cref{thm:kr-marked}, \cref{prop:kr-determinant}, \cref{lem:kr-holder}\\
\cref{thm:ent-main} & conditioned hitting estimate & \cref{lem:ent-regularity}, \cref{prop:ent-regression}, \cref{prop:ent-multipliers}, \cref{lem:ent-resolvent}, \cref{prop:ent-truncated}, \cref{lem:ent-boundary}, \cref{lem:red-kernels}\\
\cref{prop:th-structure} & static consequences of strict 1RSB & \cref{lem:pre-support}, \cref{thm:th-contact}, \cref{thm:th-subag}\\
\cref{thm:dir-norm} & directing vectors have squared norm $Q$ & \cref{prop:th-structure}, \cref{lem:th-replica-exclusion}, \cref{lem:dir-cap}, \cref{lem:dir-exclusion}, \cref{lem:dir-GG}\\
\cref{thm:mc-counts} & counts and existence, non-pure mixtures & \cref{thm:mo-microscopic}; exclusion \cref{thm:mo-filter,prop:mo-filter-loss}; pair bound \cref{thm:mp-local-pair,cor:mp-filter-upper}; \cref{lem:mp-gram,thm:mp-count}; finite counts \cref{thm:kr-marked,cor:kr-count}\\
\cref{prop:pc-coarse}, \cref{prop:mb-coarse} & coarse cap cover in pure models; coarse band cover in non-pure mixtures & \cref{thm:dir-norm}; replica means \cref{prop:pc-replica-means} or \cref{lem:mb-replica}; localization \cref{prop:pc-relocation} or \cref{lem:mb-localization}, with \cref{prop:mo-all-high,lem:mb-counts}\\
\cref{thm:mb-capture} & concentration, non-pure mixtures & \cref{prop:red-capture}, \cref{lem:mb-counts}, \cref{prop:mb-coarse}, \cref{cor:mb-bad}, \cref{cor:mb-low-sum}; measurability and finite critical sets \cref{lem:mb-borel,thm:kr-marked}\\
\cref{thm:pc-capture} & concentration, pure models & \cref{prop:pc-coarse,lem:pc-diagonal}, \cref{lem:pc-palm}, \cref{prop:pc-bands}, \cref{prop:pc-laplace}, \cref{prop:pc-denominator}, \cref{lem:pc-tight}\\
\bottomrule
\end{tabular}
\caption{Dependencies in \cref{part:condensed}.}
\label{tab:ov-condensed}
\end{table}

\begin{table}[ht]
\centering
\small
\begin{tabular}{@{}>{\raggedright\arraybackslash}p{0.21\textwidth}>{\raggedright\arraybackslash}p{0.30\textwidth}>{\raggedright\arraybackslash}p{0.41\textwidth}@{}}
\toprule
Result & Content & Uses\\
\midrule
\cref{thm:main-memory} & slow relaxation of the equilibrium dynamics & \cref{prop:pr-certificate}, \cref{prop:pr-good}, \cref{cor:pr-cap}, \cref{lem:sp-barycenter}, \cref{lem:sp-bumps}, \cref{lem:sp-projection}, \cref{lem:sp-trap}\\
\cref{prop:pr-good} & bounds on caps, core bands, and collars around good reference points & \cref{lem:pr-latitude}, \cref{prop:pr-cross}, \cref{prop:pr-certificate}, \cref{lem:pr-transfer}\\
\cref{prop:pr-certificate} & the equality at $q_*$ & \cref{lem:pr-trial}, \cref{lem:pre-support}\\
\cref{cor:pr-cap} & uniform bound on the Gibbs masses of caps & \cref{prop:pr-replica}, \cref{lem:pr-cap}\\
\cref{lem:sp-bumps} & many slow modes from one aligned observable & \cref{lem:sp-frame}\\
\cref{thm:main-plateau} & correlation bounds at every level below $q_*$ & \cref{thm:main-memory}, \cref{prop:pl-norm}, \cref{lem:pl-energy}, \cref{lem:sp-projection}, \cref{cor:pr-cap}\\
\cref{prop:pl-norm} & the norm of the overlap-weighted Gibbs mean & \cref{prop:pl-planted}, \cref{lem:pl-pressure}, \cref{lem:pl-envelope}, \cref{lem:pl-ibp}; contiguity (\cref{sec:pl-transfer})\\
\cref{thm:main-cusp} & eigenvalues near the dynamical threshold & \cref{lem:inf-field}, \cref{lem:inf-dynamics}, \cref{lem:inf-pressure}, \cref{lem:inf-good}, \cref{lem:inf-frame}, \cref{lem:inf-replica}\\
\bottomrule
\end{tabular}
\caption{Dependencies in \cref{part:prestatic}.}
\label{tab:ov-prestatic}
\end{table}
